\section{Planteamiento de la localización y sistemas estudiados}
\label{sec:master-scope}

La localización de un atractor oculto comienza con una dificultad concreta:
iniciar cerca de los equilibrios puede llevar a otros destinos. Necesitamos
construir condiciones iniciales en regiones de la cuenca que no conocemos de
antemano. Para ello combinamos información del campo, aproximaciones
armónicas y continuación. Una vez obtenida la trayectoria, estudiamos su
atracción, su régimen dinámico y su relación con las vecindades de equilibrio.

Cada operación responde a una pregunta distinta. Resolver una identidad del
campo permite transformar las ecuaciones; integrar una semilla permite
examinar su destino; sondear vecindades permite estudiar el acceso a ese
destino. Esta relación se resume en la \cref{tab:master-levels}.

\begin{longtable}{@{}L{0.23\linewidth}L{0.30\linewidth}L{0.39\linewidth}@{}}
\caption{Operaciones de localización e información que proporcionan.}
\label{tab:master-levels}\\
\toprule
Operación & Magnitud calculada & Interpretación dinámica\\\midrule
\endfirsthead
\toprule Operación & Magnitud calculada & Interpretación dinámica\\\midrule
\endhead
Realización de Lur'e & Matrices, no linealidad y resolvente & Identifica la
estructura de realimentación del campo.\\
Balance armónico & Frecuencia, amplitud y sesgo & Propone una condición
oscilatoria aproximada para iniciar la búsqueda.\\
Geometría del campo & Equilibrios, jacobiano y puntos perpetuos & Organiza
regiones y condiciones iniciales mediante propiedades locales.\\
Continuación & Trayectoria a lo largo de un parámetro auxiliar & Transporta
una respuesta del sistema auxiliar al sistema de interés.\\
Diagnósticos de trayectoria & Retornos, espectros y crecimiento variacional &
Caracteriza la respuesta en el intervalo de observación.\\
Sondeos de vecindades & Destinos por equilibrio y radio & Examina si las
condiciones próximas a los equilibrios alcanzan la respuesta localizada.\\
Comparación independiente & Diferencias entre métodos, pasos y lenguajes &
Evalúa la sensibilidad de las conclusiones a la discretización.\\
\bottomrule
\end{longtable}

\subsection{Sistemas estudiados}

Los modelos enteros permiten trabajar con un flujo sobre un espacio de
estados finito dimensional. En los modelos de Caputo fijamos además el
terminal inferior y la historia. Comparamos los casos de
\cref{tab:master-cases} bajo estas dos formulaciones, conservando los
parámetros y las condiciones iniciales de cada experimento.

\begin{longtable}{@{}L{0.25\linewidth}C{0.08\linewidth}L{0.59\linewidth}@{}}
\caption{Sistemas, órdenes y respuestas observadas.}\label{tab:master-cases}\\
\toprule
Sistema & Orden & Resultado principal\\\midrule
\endfirsthead
\toprule Sistema & Orden & Resultado principal\\\midrule
\endhead
Chua PWL canónico & \(q=1\) & Semilla de función descriptiva y continuación.
La comparación terminal de 504 condiciones distingue 168 destinos de
equilibrio, 171 cruces de norma y 165 respuestas inconclusas.\\
Chua PWL cercano & \(q=0.9998\) & Control autoexcitado: 1305 contactos de
4050 condiciones y 1395 divergencias bajo el protocolo ensayado.\\
Chua PWL de Danca, centrado & \(q=0.9998\) & Dos semillas armónicas para los
parámetros nominales, distintas de la semilla sesgada.\\
Chua PWL de Danca, sesgado & \(q=0.9998\) & Oscilación regular con frecuencia
dominante reproducida. Los sondeos locales dan \(0/7200\) contactos y la
extensión a \(r=0.3\) registra 37.\\
Chua arctangente de Wu & \(q=0.99\) & Respuestas simétricas regulares con
ADM local; este método reinicia la memoria en cada subintervalo.\\
Chua arctangente c590 & \(q=0.9999\) & Respuesta aperiódica en la
exploración inicial y sensibilidad posterior al integrador y al paso.
Los sondeos registran \(0/8400\), \(22/2400\) y \(588/2700\) contactos
en los bloques de radio hasta \(0.3\), \(1\) y \(2\), respectivamente.\\
Kalman--Fitts & \(q=1\) & Ciclo estable localizado mediante mapa de
conmutación y homotopía signo--tanh, con una obstrucción algebraica del
balance armónico directo.\\
MAVPD, \(\xi=3.1\) & \(q=1\) & Ciclo periódico externo de cuenca oculta en
la clasificación numérica: \(0/108\) contactos locales. Las semillas PP
exactas alcanzan las ramas interiores simétricas.\\
MAVPD, \(\xi=3.5\) & \(q=1\) & Control negativo con 24 contactos desde las
vecindades: la cuenca es accesible bajo el protocolo común.\\
MAVPD, \(\xi=2.85\) & \(q=1\) & Régimen caótico para
\(\gamma=0.15380379839949113\), con \(0/112\) contactos locales;
Python y Julia coinciden en las 112 decisiones de destino.\\
PLL lead--lag & \(q=1\) & Ciclo de rotación sobre el cilindro de fases y
\(0/96\) contactos locales; el cierre armónico directo está obstruido.\\
Lorenz generalizado & \(q=0.995\) & Realización MIMO, equilibrios y dos
semillas FPP--A. Las integraciones desde esas semillas se aproximan a los
equilibrios simétricos.\\
Rabinovich--Fabrikant & \(q=0.998\) & Realización MIMO, cinco equilibrios,
clasificación de Matignon y ocho semillas FPP--A.\\
Muñoz--Pacheco & \(q=0.97\) & Cuatro semillas FPP--A exactas y
trayectorias PECE de historia completa hasta \(t=35\).\\
Otras familias & Mixto & Construcciones compatibles con Lur'e y con la
búsqueda geométrica, cuyo destino requiere integración específica.\\
\bottomrule
\end{longtable}

En cada fila, un contacto es una decisión calculada con una referencia,
una distancia y una ventana temporal determinadas. La comparación entre
semillas emplea esa misma definición de destino. Las ramas y las semillas
que llegan a una misma respuesta se agrupan antes de contar respuestas
distintas; así evitamos identificar el número de condiciones iniciales con
el número de atractores.

\subsection{Convención de realimentación}

Consideramos la representación
\begin{equation}
 {}^{\mathrm C}D_{t_0}^{q}x=Px+b\,\psi(r^{\mathsf T}x),\qquad
 W_+(z)=r^{\mathsf T}(zI-P)^{-1}b,\qquad z=s^q.
 \label{eq:master-wplus}
\end{equation}
El orden se denota por \(q\) y el vector de entrada por \(b\). Si aproximamos
la no linealidad mediante una ganancia real \(N\), la salida fundamental
\(\Sigma\ne0\) satisface \(\Sigma=W_+(z)N\Sigma\). Al pasar todo al
primer miembro y dividir por \(\Sigma\), obtenemos
\begin{equation}
 1-W_+(z)N=0.
 \label{eq:master-closure}
\end{equation}
Escribiendo \(W_+=W_R+\ii W_I\), el cierre se separa en
\[
 1-NW_R=0,\qquad NW_I=0.
\]
La primera igualdad impide \(N=0\). Por tanto, necesitamos \(W_I=0\) y
\(W_R\ne0\), y la ganancia queda fijada por \(N=1/W_R\). Aún debemos
comprobar que esa ganancia pertenece al rango de la función descriptiva.

Si usamos la resolvente en el orden contrario,
\begin{equation}
 \begin{aligned}
 W_-(z)&=r^{\mathsf T}(P-zI)^{-1}b
        =-r^{\mathsf T}(zI-P)^{-1}b=-W_+(z),\\
 1+W_-(z)N&=0.
 \end{aligned}
 \label{eq:master-wminus}
\end{equation}
La igualdad \((P-zI)^{-1}=-(zI-P)^{-1}\) vale cuando la resolvente existe.
Así, \(-1/\Re W_-=1/\Re W_+\): ambas convenciones dan la misma condición
si el signo del cierre se cambia junto con el de la transferencia.

\subsection{Escalas y variables de comparación}

La semilla canónica de Chua entero tiene amplitud
\(A_0=5.856145086257\). Al introducir un sesgo distinguimos la
transferencia original de la incremental: la pendiente local se incorpora
una sola vez en el bloque lineal. En la familia auxiliar de Machado
evaluamos por separado el cierre con \(N^\mu\) y el balance físico con
\(N\), porque corresponden a ecuaciones diferentes.

La amplitud espectral, la potencia y la densidad espectral de potencia
tienen normalizaciones y unidades distintas. Del mismo modo, una semilla
FPP--A es una condición algebraica del arranque de Picard--Caputo, mientras
que un evento FPP--H se evalúa sobre una trayectoria con historia. Estas
distinciones fijan las magnitudes que podemos comparar sin cambiar de
problema dinámico.

Para clasificar la ocultedad estudiamos el acceso a la respuesta desde
vecindades de todos los equilibrios. En los sondeos especificamos radios,
distribución de condiciones iniciales, horizonte y tolerancias. Un argumento
analítico de separación debe cubrir una vecindad completa; en Caputo debe
precisar también la clase de historias admisibles. Cuando no existen
equilibrios, la condición de separación respecto de ellos es vacía, pero
la existencia y la atracción del conjunto siguen siendo propiedades que
hay que establecer.
